\section{2014 Algebra & Number Theory}

\subsection*{Individual}
\begin{exercise}
Let $G$ be a finite subgroup of $GL(V)$ where $V \cong \mathbb{C}^n$.
(a) $H = \{ h \in G : h|_V = \eta(h) I \}$ is normal and $H \cong \text{im}(\eta)$.
(b) $|\chi_V(g)| \leq n$, with equality iff $g \in H$.
(c) Irrep $W$ is a direct summand of $V^{\otimes m}$ for some $m$.
\end{exercise}

\begin{solution}
(a) $H$ consists of those elements in $G$ that act as scalars on $V$.
For any $h \in H$, $h = \eta(h) I$. For any $g \in G$, $g h g^{-1} = g (\eta(h) I) g^{-1} = \eta(h) g g^{-1} = \eta(h) I \in H$. Thus $H \triangleleft G$.
The map $\eta: H \to \mathbb{C}^\times$ is a group homomorphism. If $\eta(h) = 1$, then $h = I$, so it is injective.

(b) \textbf{Pedagogical Remark:} The character value is the sum of eigenvalues, which are roots of unity for a finite group.
Let $\lambda_1, \dots, \lambda_n$ be the eigenvalues of $g$. Since $g^k = I$ for some $k$, each $\lambda_j$ is a $k$-th root of unity, so $|\lambda_j| = 1$.
$|\chi_V(g)| = |\sum \lambda_j| \leq \sum |\lambda_j| = n$ by the triangle inequality.
Equality $|\sum \lambda_j| = \sum |\lambda_j|$ holds if and only if all $\lambda_j$ lie on the same ray from the origin. Since $|\lambda_j|=1$, this means $\lambda_1 = \dots = \lambda_n$.
Then $g$ is a diagonalizable matrix with all eigenvalues equal to some $\lambda$, so $g = \lambda I \in H$.

(c) \textbf{Theorem (Burnside):} If $V$ is a faithful representation of a finite group $G$, then every irreducible representation of $G$ occurs as a constituent of some $V^{\otimes m}$ (for $m \geq 0$).
If $G \subset GL(V)$, then $V$ is faithful by definition.
The tensor algebra $T(V) = \bigoplus_{m=0}^\infty V^{\otimes m}$ contains all irreducible representations.
This can be seen using characters: if $W$ is an irrep not in any $V^{\otimes m}$, then $\langle \chi_W, \sum \chi_V^m \rangle = 0$, which implies a contradiction with the density of polynomials in the space of functions on the group.
\end{solution}

\begin{exercise}
(a) $A = (t^{a_i+a_j})$ is PSD. Rank. (b) $B = (1/(1+a_i+a_j))$ is PSD. (c) $B$ is PD iff $a_i$ distinct.
\end{exercise}

\begin{solution}
(a) Let $x_i = t^{a_i}$. Then $A_{ij} = x_i x_j$.
For any vector $v$, $v^T A v = \sum v_i x_i x_j v_j = (\sum v_i x_i)^2 \geq 0$.
The rank is 1 as $A = \mathbf{x} \mathbf{x}^T$ for the non-zero vector $\mathbf{x} = (t^{a_1}, \dots, t^{a_n})^T$.

(b) \textbf{Motivation:} This is a \textit{Gram matrix} for a certain inner product on functions.
Using the identity $1/S = \int_0^\infty e^{-sx} dx$:
$B_{ij} = \int_0^\infty e^{-(1+a_i+a_j)x} dx = \int_0^\infty e^{-x} e^{-a_i x} e^{-a_j x} dx$.
Then $v^T B v = \int_0^\infty e^{-x} (\sum v_i e^{-a_i x})^2 dx \geq 0$.

(c) $B$ is positive definite if $v^T B v = 0 \implies v = 0$.
The integral is 0 iff the integrand $f(x) = \sum v_i e^{-a_i x}$ is 0 almost everywhere.
By analyticity, $f(x) \equiv 0$.
The functions $\{ e^{-a_i x} \}$ are linearly independent iff $a_i$ are distinct.
Thus $v = 0$ is the only solution iff $a_i$ are distinct.
\end{solution}

\begin{exercise}
$X^2 - 82Y^2 = \pm 2$.
(a) $(x, y) \to (9x-82y, x-9y)$. (b) Roots mod $p^n$. (c) No integer roots.
\end{exercise}

\begin{solution}
(a) Norm calculation: $N(x+y\sqrt{82}) = x^2 - 82y^2$.
The map corresponds to multiplication by $9-\sqrt{82}$, which has norm $81-82 = -1$.
$N((x+y\sqrt{82})(9-\sqrt{82})) = (x^2-82y^2)(-1) = \mp 2$.
Explicitly: $(9x-82y)^2 - 82(x-9y)^2 = 81x^2 - 2(9)(82)xy + 82^2y^2 - 82x^2 + 2(82)(9)xy - 82(81)y^2$
$= (81-82)x^2 + (82^2 - 82 \cdot 81)y^2 = -x^2 + 82y^2 = -(x^2-82y^2)$.

(b) By Hasse Principle, we check local solubility.
For $p$ odd, $x^2 - 82y^2 \equiv 2 \pmod p$ has a solution if $(\frac{82}{p}) = (\frac{2}{p})$ or $(\frac{-82}{p}) = (\frac{-2}{p})$.
This is always true as $82 = 2 \cdot 41$. The condition is about Hilbert symbols.
Actually, $2$ is a norm in $\mathbb{Q}_p(\sqrt{82})$ for all odd $p$.

(c) \textbf{Pell-like equations and units:}
The fundamental unit of $\mathbb{Q}(\sqrt{82})$ is $\epsilon = 9 + \sqrt{82}$ (norm -1).
The equation $x^2 - 82y^2 = 2$ means the ideal $(2)$ is a square of a principal ideal $(x+y\sqrt{82})$.
In $\mathbb{Q}(\sqrt{82})$, $2$ ramifies: $(2) = \mathfrak{p}_2^2$.
Is $\mathfrak{p}_2$ principal? The continued fraction of $\sqrt{82}$ is $[9, \overline{18}]$.
The only norms $N(x+y\sqrt{82})$ that can occur with $|N| < \sqrt{82}$ are the denominators of the continued fraction algorithm.
For $[9, \overline{18}]$, the values $Q_i$ are $Q_0=1, Q_1=1$.
Thus only norm $\pm 1$ are possible. $2$ is not a norm.
\end{solution}

\begin{exercise}
$G = S \times T$ non-abelian simple.
(a) 4 normal subgroups. (b) $S \cong T \implies \exists$ maximal not containing factors. (c) Converse.
\end{exercise}

\begin{solution}
(a) \textbf{Theorem (Goursat's Lemma):} Subgroups of a product $S \times T$ correspond to isomorphisms between quotients.
Let $N \triangleleft S \times T$. The projections $p_S(N)$ and $p_T(N)$ are normal in $S$ and $T$.
Since $S$ is simple, $p_S(N) \in \{1, S\}$.
If $N \cap (S \times 1) = S \times 1$, then $S \times 1 \subseteq N$. Then $N$ is either $S \times 1$ or $S \times T$.
If $N \cap (S \times 1) = 1$, and $p_S(N) = S$, then $N$ is a graph of an isomorphism $S \to T'$ for some $T' \triangleleft T$.
But for a graph to be normal, the isomorphism must commute with all inner automorphisms, meaning the image is in the center. Since $S, T$ are non-abelian simple, $Z(S)=1, Z(T)=1$.
Thus $N$ must be $1, S \times 1, 1 \times T, S \times T$.

(b) If $S \cong T$, let $\phi: S \to T$ be an isomorphism.
The diagonal subgroup $\Delta = \{ (s, \phi(s)) : s \in S \}$ is a subgroup of $G$.
It is maximal and contains neither $S \times 1$ nor $1 \times T$.

(c) If $M < S \times T$ is maximal and does not contain $S \times 1$, then $M \cap (S \times 1) = 1$ (by simplicity).
By Goursat, $M$ is a graph of an isomorphism between $S$ and $T$.
Thus $S \cong T$.
\end{solution}

\begin{exercise}
Chevalley-Warning Theorem.
\end{exercise}

\begin{solution}
\textbf{Goal:} $n > \sum d_i \implies N \equiv 0 \pmod p$.
Consider the polynomial $P(X) = \prod_{i=1}^r (1 - f_i(X)^{q-1})$.
$P(X) = 1$ if $X$ is a common root, and $P(X) = 0$ otherwise.
The number of roots is $N \equiv \sum_{X \in \mathbb{F}^n} P(X) \pmod p$.
The degree of $P$ is $(q-1) \sum d_i < (q-1)n$.
Any monomial in $P$ is $\prod X_j^{a_j}$ with $\sum a_j < (q-1)n$.
At least one $a_k < q-1$.
The sum $\sum_{x \in \mathbb{F}} x^a$ is $-1$ if $(q-1) \mid a, a > 0$, and 0 if $0 < a < q-1$, and $q$ if $a=0$.
In all cases for $0 \leq a < q-1$, the sum is $0$ in $\mathbb{F}$ (except for $a=0$ where it is $q \equiv 0$).
Wait, for $a=0$, $\sum 1 = q \cdot 1 = 0$.
So $\sum_{X} \prod X_j^{a_j} = \prod (\sum x_j^{a_j}) = 0$.
Thus $N \equiv 0 \pmod p$. Since $(0, \dots, 0)$ is a root, $N \geq 1$, so $N \geq p$.
\end{solution}

\begin{exercise}
(a) $AB=BA \implies \det(A^2+B^2) \geq 0$. (b) $k$ commuting matrices.
\end{exercise}

\begin{solution}
(a) $A^2 + B^2 = (A+iB)(A-iB)$.
Since $A, B$ are real, $A-iB = \overline{A+iB}$.
$\det(A^2+B^2) = \det(A+iB) \det(\overline{A+iB}) = \det(A+iB) \overline{\det(A+iB)} = | \det(A+iB) |^2 \geq 0$.

(b) For $k$ commuting matrices $A_m$, they are simultaneously triangularizable over $\mathbb{C}$.
The eigenvalues of $\sum A_m^2$ are $\sum \lambda_{m,j}^2$.
Since $A_m$ are real, the eigenvalues come in conjugate pairs? No, the \textit{vectors} of eigenvalues $(\lambda_{1,j}, \dots, \lambda_{k,j})$ come in conjugate pairs.
Contribution of a pair $(v, \bar{v})$ is $(\sum \lambda_m^2)(\sum \bar{\lambda}_m^2) = | \sum \lambda_m^2 |^2 \geq 0$.
Contribution of a real eigenvalue vector is $\sum \lambda_m^2 \geq 0$.
Product is $\geq 0$.
\end{solution}

\subsection*{Team}
\begin{exercise}
$g$ orthogonal. (1) $a=(g-1)b \neq 0 \implies \ker(s_a g - 1) = \ker(g-1) \oplus \mathbb{R}b$.
(2) $g$ is product of $k = \dim (g-1)V$ reflections.
\end{exercise}

\begin{solution}
(1) $s_a g x = x \iff g x = s_a x = x - 2\frac{(x, a)}{(a, a)} a$.
This means $(g-1)x$ is proportional to $a = (g-1)b$.
If $x \in \ker(g-1)$, then $(g-1)x = 0$ and $(x, a) = (gx-x, b) = 0$, so works.
For $x=b$, $(g-1)b = a$. We need $-2\frac{(b, a)}{(a, a)} = 1$.
$(a, a) + 2(b, a) = (gb-b, gb-b) + 2(b, gb-b) = \|gb\|^2+\|b\|^2-2(gb, b) + 2(b, gb) - 2\|b\|^2 = 0$.
So $b$ works.

(2) Induction on $k = \dim \text{im}(g-1)$.
If $k > 0$, pick $b$ such that $(g-1)b \neq 0$. Let $g' = s_a g$.
Then $\dim \ker(g'-1) = \dim \ker(g-1) + 1$, so $k' = k-1$.
By induction $g' = s_1 \dots s_{k-1}$.
Thus $g = s_a s_1 \dots s_{k-1}$ is a product of $k$ reflections.
\end{solution}

\begin{exercise}
$G$ semi-direct product $(\mathbb{Z}/p) \ltimes (\mathbb{Z}/q)^n$.
\end{exercise}

\begin{solution}
\textbf{Motivation:} This is a classification of groups with very restricted element orders.
Since $H_q \triangleleft G$ is an abelian $q$-group and all elements have order $q$, $H_q \cong \mathbb{F}_q^n$.
$G = H_q \rtimes P$. $P$ is a $p$-group with all elements order $p$, so $P \cong \mathbb{F}_p^m$.
No element has order $pq \implies P$ acts fixed-point-freely on $H_q \setminus \{0\}$.
An abelian group acting fixed-point-freely on another must be cyclic (or have very limited structure).
Thus $P \cong \mathbb{Z}/p\mathbb{Z}$.
The action $M(1)$ has no eigenvalue 1 in $\mathbb{F}_q$, so all eigenvalues are roots of $x^p-1=0$ other than 1.
\end{solution}

\begin{exercise}
$\zeta = 1 + N \eta \implies \zeta=1$ for $N \geq 3$.
\end{exercise}

\begin{solution}
\textbf{Explanatory Remark:} This is about the "size" of algebraic integers.
Let $\eta = (\zeta-1)/N$. For any conjugate $\zeta_i$ of $\zeta$:
$|\eta_i| = | \zeta_i - 1 | / N \leq 2/N \leq 2/3 < 1$.
The norm $N(\eta) = \prod \eta_i$ is an integer in $\mathbb{Z}$.
But $|N(\eta)| = \prod |\eta_i| < 1$.
Thus $N(\eta) = 0$, so $\eta = 0$ and $\zeta = 1$.
\end{solution}

\begin{exercise}
Molien's Theorem generalization.
\end{exercise}

\begin{solution}
$\sum_{n \geq 0} \chi_{S^n V}(g) t^n = \frac{1}{\det(I - \pi(g) t)}$.
The multiplicity $a_n(\rho) = \langle \chi_{S^n V}, \chi_\rho \rangle = \frac{1}{|G|} \sum_g \chi_{S^n V}(g) \overline{\chi_\rho(g)}$.
Summing over $n$:
$\sum a_n(\rho) t^n = \frac{1}{|G|} \sum_g \overline{\chi_\rho(g)} \sum_n \chi_{S^n V}(g) t^n = \frac{1}{|G|} \sum_g \frac{\overline{\chi_\rho(g)}}{\det(I - \pi(g) t)}$.
\end{solution}

\begin{exercise}
$A$ normal $\iff$ unitarily diagonalizable.
\end{exercise}

\begin{solution}
(1) $\ker(T)^\perp = \text{im}(T^*)$ is a standard result.
$\ker(A-\lambda)^\perp = \text{im}(A^*-\bar{\lambda})$.

(2) \textbf{Spectral Theorem:} A matrix $A \in M_n(\mathbb{C})$ is normal ($AA^* = A^*A$) if and only if there exists a basis of eigenvectors, which is equivalent to being unitarily diagonalizable.
$(\impliedby)$ If $A = U D U^*$, then $A^* = U D^* U^*$.
$A A^* = U D D^* U^* = U D^* D U^* = A^* A$ since diagonal matrices commute.
\end{solution}

\begin{exercise}
$x^5 - 80x + 5$ has Galois group $S_5$.
\end{exercise}

\begin{solution}
(1) Irreducible by Eisenstein at $p=5$.
(2) (a) $f'(x) = 5x^4 - 80 = 5(x^4-16)$. Roots $x = \pm 2$.
$f(2) = 32 - 160 + 5 = -123$. $f(-2) = -32 + 160 + 5 = 133$.
Sign change at $-\infty, -2, 2, \infty$ gives 3 real roots.
(b) Irreducible degree 5 $\implies$ 5-cycle. 2 complex roots $\implies$ transposition (complex conjugation).
(c) $S_5$ is generated by any 5-cycle and any transposition.
\end{solution}
